%!TEX TS-program = xelatex
%!TEX TS-program = xelatex
%!TEX encoding = UTF-8 Unicode
% One-page condensed version of cv.tex — same Awesome-CV template, trimmed content.
% Compile with: xelatex cv-onepage.tex

%-------------------------------------------------------------------------------
% CONFIGURATIONS
%-------------------------------------------------------------------------------
\documentclass[11pt, a4paper]{cv}
\usepackage{datetime}

% Tighter margins to fit one page
\geometry{left=1.2cm, top=0.6cm, right=1.2cm, bottom=1.2cm, footskip=.48cm}

\colorlet{awesome}{awesome-red}
\setbool{acvSectionColorHighlight}{true}
\renewcommand{\acvHeaderSocialSep}{\quad\textbar\quad}

% Reduce vertical spacing vs. default template
\renewcommand{\acvHeaderAfterSocialSkip}{2.5mm}
\renewcommand{\acvSectionTopSkip}{2mm}
\renewcommand{\acvSectionContentTopSkip}{1.5mm}

%-------------------------------------------------------------------------------
%	PERSONAL INFORMATION
%-------------------------------------------------------------------------------
\name{Eduardo}{Rodríguez Sánchez}
\position{Computer Engineer{\enskip\cdotp\enskip}HPC Master Student}
\address{Madrid, Spain}

\email{rod.san.eduardo@gmail.com}
\github{nekronos-gh}
\linkedin{eduardo-rodriguez-sanchez}
\homepage{https://erodriguez.dev/}
% Quote removed to save space on one-page version

%-------------------------------------------------------------------------------
\begin{document}

\makecvheader
\makecvfooter
  {\monthname[\the\month] \the\year}
  {Eduardo R. S.~~~·~~~Curriculum Vitae}
  {\thepage}

%-------------------------------------------------------------------------------
%	EDUCATION
%-------------------------------------------------------------------------------
\cvsection{Education}
\begin{cventries}
  \cventry
  {M.S. in High Performance Computing --- EUMaster4HPC double degree} % Degree
  {Univ. of Luxembourg / Sofia Univ. St. Kliment Ohridski} % Institution
  {Belval / Sofia} % Location
  {Sep. 2025 - Present} % Date(s)
  {
    \begin{cvitems}
      \item {Focus: parallel algorithms, HPC system architecture, accelerator programming, and performance analysis. \\ GPA: 17.5/20.}
    \end{cvitems}
  }

  \cventry
  {B.S. in Computer Engineering (exchange year: Georgia Tech)} % Degree
  {UC3M / Georgia Institute of Technology} % Institution
  {Madrid / Georgia, USA} % Location
  {Sep. 2020 - Jun. 2024} % Date(s)
  {
    \begin{cvitems}
      \item {Coursework: Operating Systems, Networks, Distributed Systems, Machine Learning, Cryptography. \\ GPA: 9.1/10 (UC3M) $\cdot$ 4.0/4.0 (Georgia Tech). \\ Thesis: C++ load-balancing proxy sustaining 86\,Gb/s of trace forwarding across millions of concurrent connections.}
    \end{cvitems}
  }
\end{cventries}

%-------------------------------------------------------------------------------
%	SKILLS
%-------------------------------------------------------------------------------
\cvsection{Skills}
\begin{cvskills}
  \cvskill
    {Programming}
    {C, C++, Python}
  \cvskill
    {HPC / Systems}
    {OpenMP, MPI, AVX2, Multithreading, Profiling, Linux, GDB, Valgrind}
  \cvskill
    {Distributed / Cloud}
    {Kubernetes, Spark, 5G Networking, gRPC, Prometheus, Grafana}
  \cvskill
    {Languages}
    {Spanish (native), English (fluent)}
\end{cvskills}

%-------------------------------------------------------------------------------
%	EXPERIENCE
%-------------------------------------------------------------------------------
\cvsection{Experience}
\begin{cventries}
  \cventry
    {Software Developer} % Job title
    {Ericsson} % Organization
    {Madrid, Spain} % Location
    {Dec. 2024 - Aug. 2025} % Date(s)
    {
      \begin{cvitems}
        \item {Shipped User Plane Analytics features and antenna geo-redundancy for the 5G Core, and cut multi-threaded lock contention to resolve critical production issues.}
        \item {Led the migration of the 5G RIB codebase to the Madrid site and redesigned intra-forwarding-module communication to run entirely inside Kubernetes, eliminating external network hops.}
      \end{cvitems}
    }

  \cventry
    {Software Engineer} % Job title
    {Samsung Electronics (Zhilabs)} % Organization
    {Madrid, Spain} % Location
    {Jun. 2023 - Oct. 2024} % Date(s)
    {
      \begin{cvitems}
        \item {Built a network trace proxy sustaining 86\,Gb/s with a recursive-descent binary parser for trace analysis, and parallelized a multi-antenna traffic simulator to speed up testing.}
        \item {Designed a RAN Assistant, a proof-of-concept multi-agent LLM/NLP/ML system for telecom Q\&A, anomaly detection, and automated root-cause analysis.}
      \end{cvitems}
    }
\end{cventries}

%-------------------------------------------------------------------------------
%	RESEARCH
%-------------------------------------------------------------------------------
\cvsection{Research}
\begin{cventries}
  \cventry
    {Research Assistant}
    {SnT (Security, Reliability and Trust)}
    {Luxembourg}
    {Oct. 2025 - Apr. 2026}
    {
      \begin{cvitems}
        \item {Prototyped a Byzantine fault-tolerant PBFT consensus engine in C++ on an async networking stack, including a cross-replica-group coordination protocol.}
        \item {Built failure-injection scenarios and reproducible benchmarking tooling to validate consensus performance under adversarial conditions.}
      \end{cvitems}
    }
\end{cventries}

%-------------------------------------------------------------------------------
%	PROJECTS
%-------------------------------------------------------------------------------
\cvsection{Projects}
\begin{cventries}
  \cventry
    {\href{https://github.com/nekronos-gh/rule110}{github.com/nekronos-gh/rule110}}
    {Rule110 Fast}{}{}{Cellular-automaton simulator vectorized with AVX2 and parallelized with OpenMP, using aligned memory in modern C++23.}

  \cventry
    {\href{https://github.com/nekronos-gh/TinyFile}{github.com/nekronos-gh/TinyFile} }
    {TinyFile}{}{}{Zero-copy file-compression service built on POSIX message queues and shared memory, driven by a ring-buffer state machine.}

  \cventry
    {\href{https://github.com/nekronos-gh/pbft-cpp}{github.com/nekronos-gh/pbft-cpp} }
    {PBFT}{}{}{Full 3-phase consensus implementation in C++ (Salticidae) with view changes and checkpointing, instrumented via Prometheus/Grafana.}
\end{cventries}

%-------------------------------------------------------------------------------
\end{document}
